\documentclass{report}
\usepackage{amsmath,amssymb}
\setlength{\parindent}{0mm}
\setlength{\parskip}{1em}
\begin{document}
\begin{center}
\rule{6in}{1pt} \
{\large Guangye Chen \\
{\bf An efficient mixed-precision CPU-GPU implementation of energy- and charge-conserving, fully implicit particle-in-cell algorithm}}

Fusion Energy Division \\ Oak Ridge National Laboratory \\ P O Box 2008 \\ Oak Ridge \\ TN 37831-6169
\\
{\tt gcc@ornl.gov}\\
Guangye Chen\\
Luis Chacon\\
Daniel Barnes\end{center}

A recent proof-of-principle study proposes an energy- and
charge-conserving, fully implicit particle-in-cell (PIC) algorithm for
multi-scale, full-f kinetic, electrostatic simulations in one dimension
[1]. The method employs a Jacobian-free Newton-Krylov (JFNK) solver and
is capable of taking very large time steps without loss of numerical
stability or accuracy. This is in contrast with the explicit time
integration approach (leap-frog) employed by classical PIC algorithms,
which feature stringent temporal and spatial resolution constraints to
avoid numerical instabilities [2]. By enforcing nonlinear convergence of
the field-particle equations and ensuring energy and charge conservation
to numerical round-off, the fully implicit PIC algorithm improves on
earlier implicit PIC approaches (e.g., implicit moment [3,4] and direct
implicit [5,6] methods), which failed to converge the nonlinear residual
and as a result developed large numerical errors in long-time simulations
[7], with significant self-heating/cooling often observed.

A fundamental feature of the method is the nonlinear elimination of
particle equations, which employs the enslavement of particle orbits to
the field equations. As a result of particle enslavement, it segregates
particle orbit integrations from the field solver, while remaining fully
self-consistent. This provides great flexibility, and dramatically
improves the solver efficiency by decreasing the number of degrees of
freedom of the associated nonlinear system. However, it requires a
particle push per nonlinear residual evaluation, which makes the particle
push the most time-consuming operation in the algorithm. This renders the
implicit algorithm ideally suited for emerging heterogeneous
architectures which combine CPUs and GPUs [8]. Thus, the field JFNK
solver can remain on the CPU, while the particle mover (which is
naturally data-parallel and compute-bounded) stays on the GPU. Although
the adaptivity of the particle mover creates load imbalances and dynamic
control flows (which pose a challenge for the GPU), we demonstrate that a
highly efficient GPU implementation of the implicit particle mover (using
CUDA) is possible [8]. We obtain 300 to 400 GOp/s (counting
floating-point, integer and special function operations) using single
precision, depending on the electrostatic PIC model of choice (Ampere vs.
Poisson). This is about 20% to 25% of the peak performance of the GPU
(Nvidia GeForce GTX580), and about 200 to 300 times faster than a
single-CPU (Intel Xeon@3.16GHz) serial implementation.

We have implemented an efficient mixed-precision, hybrid CPU-GPU
implementation of the implicit PIC algorithm with a defect-correction
approach [8]. In this version of the hybrid algorithm, we keep the JFNK
solver on the CPU in double precision, while the particle mover is
implemented on the GPU in both double precision (to evaluate the Jacobian
right hand side in Newton's method) and single precision (to evaluate the
Jacobian-free matrix-vector products). As a result, many single precision
particle-push evaluations are used to accelerate convergence to a
double-precision result. In a challenging long-timescale ion acoustic
wave simulation, the mixed-precision hybrid CPU-GPU solver is shown to
over-perform the double-precision CPU-only serial version by a factor of
\sim 100, without apparent loss of robustness or accuracy.

[1] G. Chen, L. Chac�n, D. C. Barnes, J. Comput. Phys. 230, (2011).
[2] C. K. Birdsall, A. B. Langdon, Plasma Physics via Computer
Simulation, McGraw-Hill, New York, 2004.
[3] R. J. Mason, J. Comput. Phys., 41 (2) (1981) 233�244.
[4] J. U. Brackbill, D. W. Forslund, Multiple time scales, Academic Press, 1985.
[5] A. Friedman, A. B. Langdon, B. I. Cohen, Plasma Phys. Controlled
Fusion, 6 (6) (1981) 225 � 36.
[6] A. B. Langdon, D. C. Barnes, Multiple time scales, Academic Press, 1985.
[7] B. I. Cohen, A. B. Langdon, D. W. Hewett, R. J. Procassini, J.
Comput. Phys., 81 (1) (1989) 151�168.
[8] G. Chen, L. Chac�n, D. C. Barnes, J. Comput. Phys. submitted (2011).


\end{document}
